% !TEX root =./ECE444_Practice_main.tex
%
% L4 practice set (Impedance, Feeding, Baluns), ported from the Jupyter Book
% practice.md / practice-solutions.md into the exam class for Gradescope.

\fullwidth{\textbf{Learning Objective 1.4: } I can calculate input impedance,
feed considerations, and the role of baluns in an antenna feed system.
\\[4pt]\normalfont\small Show your work. A \textbf{Documentation} line at the end
is required (write \textbf{None} if you did not collaborate).}

\begin{questions}

%%%%%%%%%%%%%%%%%%%%%%%%%%%%%%%%%%%%%%%%%
\question \textbf{Radiation resistance and efficiency.} A short dipole has length
$\ell = 0.08\lambda$. Its conductors contribute a loss resistance of
$R_\text{loss} = 1.5\ \ohms$.
\begin{parts}

	\begin{minipage}[t]{0.58\linewidth}
		\part Compute the radiation resistance $R_\text{rad} = 80\pi^{2}(\ell/\lambda)^{2}$.
		\begin{solution}[1in]
		\eq{ R_\text{rad} = 80\pi^{2}(0.08)^{2} = 789.6 \times 0.0064 \approx 5.05\ \ohms }
		\end{solution}
	\end{minipage}\hfill%
	\begin{minipage}[t]{0.37\linewidth}
		\raggedleft \vspace{12pt}
		$R_\text{rad}$ = \ansbox[1.3in]{$\approx 5.05\ \ohms$}
	\end{minipage}

	\begin{minipage}[t]{0.58\linewidth}
		\part Compute the radiation efficiency $\eta_\text{rad}$.
		\begin{solution}[1in]
		\eq{ \eta_\text{rad} = \dfrac{R_\text{rad}}{R_\text{rad}+R_\text{loss}} = \dfrac{5.05}{6.55} \approx 0.77 }
		\end{solution}
	\end{minipage}\hfill%
	\begin{minipage}[t]{0.37\linewidth}
		\raggedleft \vspace{12pt}
		$\eta_\text{rad}$ = \ansbox[1.3in]{$\approx 0.77$ (77\%)}
	\end{minipage}

	\begin{minipage}[t]{0.58\linewidth}
		\part The short dipole's directivity is $D = 1.5$ (1.76 dBi). Find the gain in dBi.
		\begin{solution}[1in]
		\eq{ G_\text{dBi} = 1.76 + 10\mylog{0.77} = 1.76 - 1.13 \approx 0.6\dbi }
		\end{solution}
	\end{minipage}\hfill%
	\begin{minipage}[t]{0.37\linewidth}
		\raggedleft \vspace{12pt}
		$G$ = \ansbox[1.3in]{$\approx 0.6\dbi$}
	\end{minipage}

\end{parts}

\newpage
%%%%%%%%%%%%%%%%%%%%%%%%%%%%%%%%%%%%%%%%%
\question \textbf{Feeding a half-wave dipole with $50\ \ohms$ coax.}
\begin{parts}

	\begin{minipage}[t]{0.58\linewidth}
		\part A dipole trimmed to resonance presents $\myz{in} = 70 + j0\ \ohms$. Find
		$\Gamma$, VSWR, return loss, and mismatch loss on a $50\ \ohms$ line.
		\begin{solution}[1.8in]
		\eq{
			\Gamma &= \dfrac{70-50}{70+50} = 0.167 \\
			\text{VSWR} &= \dfrac{1.167}{0.833} \approx 1.40 \\
			\text{RL} &= -20\mylog{0.167} \approx 15.6\db \\
			L_\text{mm} &= -10\mylog{1-0.167^{2}} \approx 0.12\db
		}
		\end{solution}
	\end{minipage}\hfill%
	\begin{minipage}[t]{0.37\linewidth}
		\raggedleft \vspace{12pt}
		$\Gamma$ = \ansbox[1.3in]{$0.167$ (unitless)}

		\vspace{4pt}
		VSWR = \ansbox[1.3in]{$\approx 1.40$}

		\vspace{4pt}
		RL = \ansbox[1.3in]{$\approx 15.6\db$}

		\vspace{4pt}
		$L_\text{mm}$ = \ansbox[1.3in]{$\approx 0.12\db$}
	\end{minipage}

	\begin{minipage}[t]{0.58\linewidth}
		\part An \emph{untrimmed} half-wave dipole presents
		$\myz{in} = 73 + j42.5\ \ohms$. Find $\left| \Gamma \right|$ and VSWR on the
		same line.
		\begin{solution}[1.7in]
		\eq{
			\left| \Gamma \right| &= \dfrac{\sqrt{23^{2} + 42.5^{2}}}{\sqrt{123^{2} + 42.5^{2}}} = \dfrac{48.3}{130.1} \approx 0.371 \\
			\text{VSWR} &= \dfrac{1.371}{0.629} \approx 2.18
		}
		\end{solution}
	\end{minipage}\hfill%
	\begin{minipage}[t]{0.37\linewidth}
		\raggedleft \vspace{12pt}
		$\left| \Gamma \right|$ = \ansbox[1.3in]{$\approx 0.371$}

		\vspace{4pt}
		VSWR = \ansbox[1.3in]{$\approx 2.18$}
	\end{minipage}

	\part Compare (a) and (b): in one sentence, what did trimming to resonance buy
	you, and why?
		\begin{solution}[1in]
		Trimming killed the $+j42.5\ \ohms$ of reactance, dropping VSWR from
		$\approx 2.2$ to $1.4$ --- the reactance, not the resistance, was doing most of
		the damage.
		\end{solution}

	\begin{minipage}[t]{0.58\linewidth}
		\part A third dipole on the same $50\ \ohms$ line measures VSWR $= 1.5$, and
		you know it is purely resistive at that frequency. What two real $\myz{in}$ are
		consistent with that reading? State what a VSWR measurement alone can and
		cannot tell you.
		\begin{solution}[1.5in]
		\eq{ \myz{in} = 50 \cdot \dfrac{1.5}{1} = 75\ \ohms \quad\text{or}\quad
			 \myz{in} = \dfrac{50}{1.5} = 33.3\ \ohms }
		VSWR is built from $\left| \Gamma \right|$ alone, so it reports the
		\emph{magnitude} of the mismatch and nothing about its phase --- it cannot tell
		you which side of $Z_0$ you are on, or how much of the mismatch is reactive.
		That is what a Smith chart, or a complex $\Gamma$ measurement, is for.
		\end{solution}
	\end{minipage}\hfill%
	\begin{minipage}[t]{0.37\linewidth}
		\raggedleft \vspace{12pt}
		$\myz{in}$ = \ansbox[1.5in]{$75$ or $33.3\ \ohms$}
	\end{minipage}

\end{parts}

\newpage
%%%%%%%%%%%%%%%%%%%%%%%%%%%%%%%%%%%%%%%%%
\question \textbf{Quarter-wave transformer.} A resonant antenna presents a real
input impedance $R_L = 36\ \ohms$. You want to match it to a $50\ \ohms$ feed line
with a quarter-wave transformer.
\begin{parts}

	\begin{minipage}[t]{0.58\linewidth}
		\part What characteristic impedance $Z_1$ must the $\lambda/4$ section have?
		\begin{solution}[1in]
		\eq{ Z_1 = \sqrt{Z_0 R_L} = \sqrt{(50)(36)} = \sqrt{1800} \approx 42.4\ \ohms }
		\end{solution}
	\end{minipage}\hfill%
	\begin{minipage}[t]{0.37\linewidth}
		\raggedleft \vspace{12pt}
		$Z_1$ = \ansbox[1.3in]{$\approx 42.4\ \ohms$}
	\end{minipage}

	\part What is $\Gamma$ at the design frequency?
		\begin{solution}[0.8in]
		The transformer makes $Z_1^{2}/R_L = 1800/36 = 50\ \ohms = Z_0$, so
		$\Gamma = 0$.
		\end{solution}

	\part In one sentence, what happens to the match off the design frequency, and why?
		\begin{solution}[1in]
		The section is $\lambda/4$ only at the design frequency; off frequency its
		electrical length is no longer $90\mydeg$, the transformation $Z_1^{2}/R_L$ no
		longer holds, and $\left| \Gamma \right|$ climbs --- the match is narrowband.
		\end{solution}

	\part The antenna in this problem is purely resistive. Explain in one or two
	sentences why a quarter-wave transformer is the wrong tool the moment
	$\myz{in}$ has a reactive part.
		\begin{solution}[1.2in]
		A $\lambda/4$ section transforms $R_L$ to $Z_1^{2}/R_L$ --- a \emph{real}
		number into a \emph{real} number. Hand it a complex load and the transformed
		impedance is still complex, so $\Gamma \ne 0$ no matter what $Z_1$ you pick.
		You must first cancel the reactance (or move to a matching topology that
		handles both parts at once, which is Question 4).
		\end{solution}

\end{parts}

\newpage
%%%%%%%%%%%%%%%%%%%%%%%%%%%%%%%%%%%%%%%%%
\question \textbf{L-match design.} A compact antenna presents
$\myz{in} = 25 - j40\ \ohms$ at $f = 1.5\ghz$. You must match it to a
$50\ \ohms$ line with lumped elements. Work the match in two moves: \emph{cancel
the reactance}, then \emph{transform the resistance}.
\begin{parts}

	\begin{minipage}[t]{0.58\linewidth}
		\part Find $\left| \Gamma \right|$ and VSWR for the raw, unmatched antenna.
		\begin{solution}[1.6in]
		\eq{
			\Gamma &= \dfrac{(25 - j40) - 50}{(25 - j40) + 50} = \dfrac{-25 - j40}{75 - j40} \\
			\left| \Gamma \right| &= \dfrac{\sqrt{25^{2}+40^{2}}}{\sqrt{75^{2}+40^{2}}}
				= \dfrac{47.2}{85.0} \approx 0.555 \\
			\text{VSWR} &= \dfrac{1.555}{0.445} \approx 3.49
		}
		\end{solution}
	\end{minipage}\hfill%
	\begin{minipage}[t]{0.37\linewidth}
		\raggedleft \vspace{12pt}
		$\left| \Gamma \right|$ = \ansbox[1.3in]{$\approx 0.555$}

		\vspace{4pt}
		VSWR = \ansbox[1.3in]{$\approx 3.49$}
	\end{minipage}

	\begin{minipage}[t]{0.58\linewidth}
		\part \textbf{Move 1.} What series reactance cancels the load reactance, and
		what component value does it correspond to at $1.5\ghz$?
		\begin{solution}[1.5in]
		Cancel $-j40\ \ohms$ with $+j40\ \ohms$ in series --- an inductor.
		\eq{
			\omega &= 2\pi(1.5\ee{9}) = 9.425\ee{9} \text{ rad/s} \\
			L_1 &= \dfrac{40}{9.425\ee{9}} = 4.24\ \text{nH}
		}
		The load now looks like a plain $25\ \ohms$ resistor.
		\end{solution}
	\end{minipage}\hfill%
	\begin{minipage}[t]{0.37\linewidth}
		\raggedleft \vspace{12pt}
		$X_1$ = \ansbox[1.3in]{$+j40\ \ohms$}

		\vspace{4pt}
		$L_1$ = \ansbox[1.3in]{$4.24\ \text{nH}$}
	\end{minipage}

	\begin{minipage}[t]{0.58\linewidth}
		\part \textbf{Move 2.} With $R_L = 25\ \ohms$ real, design the L-network that
		transforms it to $50\ \ohms$. Use $Q = \sqrt{Z_0/R_L - 1}$, series
		$X_s = QR_L$ on the load side and shunt $X_p = Z_0/Q$ on the line side. Give
		both reactances and both component values, and combine the two series
		inductors.
		\begin{solution}[2.0in]
		\eq{
			Q &= \sqrt{\dfrac{50}{25} - 1} = 1 \\
			X_s &= QR_L = 25\ \ohms \ (\text{series } L) \\
			L_2 &= \dfrac{25}{9.425\ee{9}} = 2.65\ \text{nH} \\
			X_p &= \dfrac{Z_0}{Q} = 50\ \ohms \ (\text{shunt } C) \\
			C &= \dfrac{1}{\omega X_p} = \dfrac{1}{(9.425\ee{9})(50)} = 2.12\ \text{pF}
		}
		The two series inductors merge: $L_1 + L_2 = 6.9\ \text{nH}$ ($+j65\ \ohms$).
		Check: $25 - j40 + j65 = 25 + j25$, whose admittance is
		$0.02 - j0.02$ S; the shunt capacitor adds $+j0.02$ S, leaving $0.02$ S
		$= 50\ \ohms$ and $\Gamma = 0$.
		\end{solution}
	\end{minipage}\hfill%
	\begin{minipage}[t]{0.37\linewidth}
		\raggedleft \vspace{12pt}
		$L_\text{tot}$ = \ansbox[1.3in]{$6.9\ \text{nH}$}

		\vspace{4pt}
		$C$ = \ansbox[1.3in]{$2.12\ \text{pF}$}
	\end{minipage}

	\part The loaded $Q$ of this match came out at $1$. Compare the bandwidth and the
	practicality of this L-match against the $\lambda/4$ transformer of Question 3,
	and say which you would build for a $1.5\ghz$ handheld radio.
		\begin{solution}[1.6in]
		A matching network's fractional bandwidth goes roughly as $1/Q$, so $Q = 1$ is
		about as broadband as a two-element match gets --- a $2{:}1$ impedance ratio is
		a gentle transformation. A larger ratio would force $Q$ up and the bandwidth
		down; that is the fundamental cost of matching, not a property of this
		topology. Against the $\lambda/4$ transformer: the transformer needs a
		\emph{real} load and a physical quarter wavelength of line ($5\ \text{cm}$ at
		$1.5\ghz$), while the L-match swallows the reactance directly and occupies two
		0402 parts. \textbf{Build the L-match} in a handheld --- board area and the
		complex load both argue for it; save the $\lambda/4$ transformer for a printed
		feed where the line is free anyway.
		\end{solution}

\end{parts}

\newpage
%%%%%%%%%%%%%%%%%%%%%%%%%%%%%%%%%%%%%%%%%
\question \textbf{Mismatch that kills the amplifier.} A power amplifier (PA)
delivers a forward power of $\myp{fwd} = +40\dbm$ toward a transmit antenna. A
shunt protection diode at the PA output fails catastrophically if it must
dissipate more than $\myp{max} = +30\dbm$. Model \emph{all} of the reflected
power as landing on the diode, so $\myp{ref} = \left| \Gamma \right|^{2}\myp{fwd}$.
\begin{parts}

	\begin{minipage}[t]{0.58\linewidth}
		\part Convert $\myp{fwd}$ and $\myp{max}$ to watts.
		\begin{solution}[1in]
		\eq{ \myp{fwd} &= 10^{(40-30)/10}\watts = 10\watts \\
			 \myp{max} &= 10^{(30-30)/10}\watts = 1\watts }
		\end{solution}
	\end{minipage}\hfill%
	\begin{minipage}[t]{0.37\linewidth}
		\raggedleft \vspace{12pt}
		$\myp{fwd}$ = \ansbox[1.3in]{$10\watts$}

		\vspace{4pt}
		$\myp{max}$ = \ansbox[1.3in]{$1\watts$}
	\end{minipage}

	\begin{minipage}[t]{0.58\linewidth}
		\part Find the largest $\left| \Gamma \right|$ the system can tolerate before
		the diode is destroyed.
		\begin{solution}[1.1in]
		\eq{ \left| \Gamma \right|_{\max} = \sqrt{\dfrac{\myp{max}}{\myp{fwd}}}
			 = \sqrt{0.1} \approx 0.316 }
		\end{solution}
	\end{minipage}\hfill%
	\begin{minipage}[t]{0.37\linewidth}
		\raggedleft \vspace{12pt}
		$\left| \Gamma \right|_{\max}$ = \ansbox[1.3in]{$\approx 0.316$}
	\end{minipage}

	\begin{minipage}[t]{0.58\linewidth}
		\part Convert to a VSWR limit and a return-loss limit, and state the
		datasheet rule that follows.
		\begin{solution}[1.5in]
		\eq{ \text{VSWR}_{\max} &= \dfrac{1.316}{0.684} \approx 1.93 \\
			 \text{RL}_{\min} &= -10\mylog{0.1} = 10\db }
		Rule: keep VSWR $\le 1.9$ (return loss $\ge 10\db$); reflected power must stay
		at least $10\db$ below forward power.
		\end{solution}
	\end{minipage}\hfill%
	\begin{minipage}[t]{0.37\linewidth}
		\raggedleft \vspace{12pt}
		VSWR$_{\max}$ = \ansbox[1.3in]{$\approx 1.93$}

		\vspace{4pt}
		RL$_{\min}$ = \ansbox[1.3in]{$10\db$}
	\end{minipage}

	\part In one sentence: how does a circulator or isolator between the PA and
	antenna change this picture?
		\begin{solution}[1in]
		It routes the reflected wave into a matched dump-port load instead of back
		into the PA, so the diode no longer sees it --- the antenna can then run at
		much higher VSWR (paid for by insertion loss, size, and bandwidth).
		\end{solution}

\end{parts}

\newpage
%%%%%%%%%%%%%%%%%%%%%%%%%%%%%%%%%%%%%%%%%
\question \textbf{Baluns.}
\begin{parts}
	\part In one or two sentences, explain why feeding a dipole \emph{directly} with
	coax (no balun) distorts the pattern.
		\begin{solution}[1.1in]
		Coax is unbalanced, so the two arms present unequal impedances to it. The
		imbalance drives a common-mode current on the \emph{outside} of the shield; that
		shield current radiates, so the feed line becomes part of the antenna and skews
		the pattern.
		\end{solution}

	\part You must feed a folded dipole ($\myz{in} \approx 292\ \ohms$) with
	$75\ \ohms$ coax. Which balun would you choose, and what does it do?
		\begin{solution}[1.1in]
		A \textbf{4:1 (half-wave) balun}: it balances the feed \emph{and} transforms
		impedance 4:1, taking $292\ \ohms$ down to $\approx 73\ \ohms$ --- a good match
		to $75\ \ohms$ coax.
		\end{solution}

	\part You feed a straight resonant dipole ($\approx 70\ \ohms$) with $50\ \ohms$
	coax. Which balun would you choose, and why is it a \emph{different} choice than (b)?
		\begin{solution}[1.2in]
		A \textbf{1:1 current (choke) balun}. The $70\ \ohms$ dipole already matches
		$50\ \ohms$ well (VSWR $\approx 1.4$), so no impedance transformation is wanted ---
		you only need to choke off the shield current. A 4:1 balun here would ruin the
		match.
		\end{solution}
\end{parts}

\vspace{1.5em}
\noindent\textbf{Documentation:} \rule[-2pt]{0.6\linewidth}{0.4pt}

\end{questions}
